\section{Discussion}
\label{sec:discussion}

\subsection{Limitations and Future Work}
\label{sec:limitations}

\paragraph{Limitations of the data sources}
While this study provides a comprehensive 13-year analysis of IDT, several limitations must be acknowledged. First, our IDT data (from the ITS survey) relies on self-reported survey responses, which may be subject to recall bias or a lack of technical knowledge regarding how information was obtained. As noted in Table \ref{tab:theft_summary_avg}, nearly 74\% of victims were unaware of the specific method used to compromise their data, limiting our ability to definitively attribute individual thefts to data breaches. Second, the structural redesign of the NCVS in 2016 introduced sampling inconsistencies as mentioned earlier: the 41\% increase in sample size and the recalibration of weights based on the 2010 Census instead of the 2000 Census could have contributed to the drop in observed OOP losses across all categories of theft. While we harmonized variables to the best of our ability, these methodological survey shifts limit our longitudinal comparisons between pre- and post-2016 survey waves. 
%
The IDT data also does not contain geographical information to allow us to account for regional price variations; as a result, our social cost model utilized fixed national averages to monetize professional services and healthcare.
%While these provide a consistent baseline, they do not account for regional price variations or individual insurance coverage, which may influence the true financial burden for specific victims. 
%
Perhaps the biggest limitations of the PRC data is the significant volume of  incidents with zero or undisclosed record exposure, which necessitated the augmentation with the HHS reports and state-level filings, a process that is inherently noisy.

%Our methodology used a conservative baseline of 1.53 million undisclosed records annually, which likely represents a significant undercounting of the true magnitude. Our reliance on linear regression to extrapolate national exposure from state level data in Maine and New Hampshire also introduces potential bias, as it assumes that breach impacts on these specific jurisdictions are representative of national trends. 

\paragraph{Limitations of our methodology}
Our attribution of IDT to specific mega-breaches relies on statistical correlation (Wilcoxon tests) rather than direct causal tracing. While we identify significant increases in reported IDTs following a mega-breach, the number of victims we attribute to major  breach events represent a modeled potential outcome.
Furthermore, our long-term social cost projections are naturally constrained by the time boundaries of the study. Because our calculations end at December 2021, the upper-bound estimates for more recent events, such as the 2017 Equifax breach, only reflect the realized portion of social cost within the study period. The true long-term cost likely extends well beyond this terminal month. As more waves of the ITS survey are released, this work can be extended to include the new data. \rev{Furthermore, as mentioned earlier, our social cost model is narrowly defined, constrained by what is captured in the ITS surveys: it focuses exclusively on harms resulting from identity theft and subsequent financial or psychological distress. Consequently, our estimates do not account for other dimensions of externalized harm, such as broader privacy costs, risks to national security, or the downstream impact of data being used for non-IDT cybercriminal activities.}

%Additionally, the breach-to-victim conversion rate is subject to the overlap or double-counting issue, where individuals are likely compromised in several breaches. Consequently, our conversion rates should be interpreted as upper bound estimates rather than precise causal links from a data breach to an IDT.

\paragraph{Future directions}
While our log-quadratic conversion model provides a robust fit for the conversion rate, it utilizes a uniform discount factor $\alpha$ to represent the decay of record utility. In reality, the ``shelf-life'' of stolen data likely varies based on the sensitivity of the identifiers involved. For example, Social Security numbers may maintain their exploitative value longer than revocable credentials like credit card numbers. Modeling this serves as a potential future extension of this study. 
%
Another line of future work involves improving the imputation process for breaches with undisclosed record counts. An alternative to our procedure (detailed in Section \ref{sec:imputation_of_undisclosed}) is to calculate a category-specific annual baseline $n_u^j$ for each of the $k$ breach types. Under this more granular approach, the imputed records for an incident $i$ of type $j$ in year $y$ would be defined as $R_{i_j}=n_u^j/|I_{u,y}^j|$. While this would account for the different typical magnitudes associated with various breach types (e.g. malicious attacks versus accidental disclosures), we maintained the global $n_u$ approach to provide a stable, conservative baseline for the overall number of undisclosed records. 
%
%more future work includes refining the social cost model to include regional cost of living adjustments and insurance policies. 
More effort can be made to identify and integrate more diverse datasets beyond the ITS and PRC datasets to better quantify the number of exposed records and number of IDTs. Finally, with the rise of generative AI in phishing and impersonation scams, research is needed to quantify the costs associated with deepfake-related IDT.
%Due to the fact that individuals could likely be compromised in multiple breaches, and our reliance on interpolated data, our calculated social costs and conversion rates should we interpreted as estimated upper bounds of the potential harm, rather than a ground truth of the specific damage caused by a single event.

\subsection{A Dynamic Saturation Model to Estimate the Number of Unique Individuals Exposed in Breaches}
\label{sec:estimating_number_records}

The analysis in Section \ref{sec:breach_social_cost} shows a potential saturation of the stolen records market, whereby even as many more records are compromised, the number of new incidents of IDT grows much slower. Below we attempt to explicitly model this saturation through a de-duplication of the record exposure.

%The analysis we present in Section \ref{sec:breach_to_victim} relied on the  augmented PRC data to identify mega-breach events and analyze the breach-to-victim conversion. While this approach provides a direct empirical baseline, we acknowledge that raw record counts likely contain significant redundancies, as an individual's personal information may be exposed across multiple independent breach events over a 13-year period. Consequently, treating every reported record as unique potentially overstates the true breadth of the victim pool. To address this, we present a dynamic saturation model as an alternative framework for de-duplicating record exposure. We acknowledge that this model is a simplification of of a complex data system. However, we offer it here as a conceptual alternative to consider when evaluating the long-term impact of data exposure.

%\subsubsection{The Dynamic Saturation Model}
The objective of this model is to transform the monthly raw number of exposed records into an estimated number of {\em unique individuals whose data was compromised}. Specifically, this model incorporates the U.S. population for ages 16+ (to match the ITS survey, which was administered to those 16+) and calculates a dynamic saturation index $\mu_t$ that tracks the portion of the population already compromised, beginning at $t=0$ and asymptotically approaching 1 as the stolen data market reaches full saturation. This ensures that as the cumulative volume of data exposure increases, the model identifies the diminishing probability that a reported record belongs to an individual whose data had not previously been compromised. 

Our model iterates using the following variables for the $t$-th period, which could be a month or a year: %\com{I suppose we can also present this on a monthly basis? we have used $t$ as the month index } \resp{Yes, I think we can as well!}


\begin{enumerate}
    \item $N_t$ (U.S. Population 16+): The potential victim pool at time $t$ based on annual population averages from the U.S. Bureau of Labor Statistics \cite{FRED_Population}.
    
    \item $r_t$ (Raw Number of Records Exposed): We take the total volume of records exposed in period $t$, $r_t$, directly from the augmented PRC data, and scale it to $\theta r_t$ with $\theta = 1.75$, to account for under-reporting as described in \cite{Graves2018Should}.  Here, we use ``records'' and not ``individuals,'' because the input to the model is number of records, and the output is unique number of individuals. 

       % For mega-breaches we \rev{could apply} an additional overlap discount of \rev{$\theta'=0.8$} to the reported record count $r_t$, so that the number of records from such an event is taken to be $\theta' r_t$ before being added to the yearly total, to model the high probability that these massive data breaches contain redundant information for individuals already present in the national supply.
       
    \item $c_t$ (Estimated Number of Newly Exposed Unique Individuals): The estimated number of previously uncompromised individuals exposed during period $t$. 
        
    \item $C_t = \sum_{i=1}^{t} c_i$ (Estimated Number of Cumulative Exposed Unique Individuals): The total volume of estimated unique individuals at the end of period $t$. 
    
    \item $\mu_t$ (Saturation Index): The ratio of the number of cumulative unique individuals at the end of period $t$ over the potential victim pool. 
    
    \item $\gamma_t$ (Dynamic Identity Density): The probability that a record in a new breach represents a newly unique individual, where $\gamma_t = \gamma_0 \times (1-\mu_{t-1})$. We set the base density as $\gamma_0 = 0.8$.
\end{enumerate}

We initialize the model at $t=0$ (corresponding to the start of 2008), where $C_0 = 0$ and $\mu_0 = 0$. Our model updates $\mu_t$ recursively. We define the update as follows, where $Y$ is a scaling factor: 

\begin{equation}
\label{eq:sat_model}
    1 - \mu_t = (1 - \mu_{t-1}) \cdot \exp\left( - \frac{\theta r_t \cdot \gamma_0 }{N_t \cdot Y } \right)~. 
\end{equation}

The role of $Y$ as a scaling factor is to define the total capacity of the identity market, effectively representing how many distinct ``identity units'' (such as different accounts) exist per person within the population. When $Y=1$, the model assumes that each individual in the population $N_t$ 
possesses exactly one ``identity unit'' that can be compromised. When $Y>1$, we scale the potential record pool of our model to $N_t \cdot Y$. This follows the assumption that one individual may have multiple distinct accounts (e.g. financial, social media, medical).

Rearranging the terms in Equation \ref{eq:sat_model} allows us to solve for the new saturation level $\mu_t$ based on the previous year's level $\mu_{t-1}$ and the influx of newly breached records $r_t$: 


\begin{equation}
\mu_t = 1 - \left[ (1 - \mu_{t-1}) \cdot \exp\left( -\frac{\theta r_t \cdot \gamma_0}{N_t \cdot Y } \right) \right]
\end{equation}


This formulation ensures that even as $r_t$ grows arbitrarily large, $\mu_t$ asymptotically approaches 1 but never exceeds it. Once the new saturation level $\mu_t$ is determined, we calculate the number of newly unique individuals $c_t$ exposed during the period by taking the difference in the cumulative pool: $$c_t = (\mu_t - \mu_{t-1}) \cdot (N_t \cdot Y)$$ This value of $c_t$ represents the net number of new individuals that were compromised during this period. 

%\subsubsection{Results from the Dynamic Saturation Model}
In Appendix \ref{app:sat_model}, we recreate Figures \ref{fig:section7_fig5}, \ref{fig:conversion_rate}, and \ref{fig:wilcoxon_raw} using the output of the above saturation model (estimated number of unique individuals compromised) rather than the number of records exposed. The implementation of the dynamic saturation model provides several key advantages for understanding the long-term relationship between data breaches and victimization by accounting for the diminishing marginal impact of redundant record exposure. As illustrated in the updated overlay analysis (Figure \ref{fig:saturation_fig5}), the estimated number of unique compromised individuals aligns much more closely with the trajectory of reported successful identity theft victims throughout the study period than the raw counts alone in Figure \ref{fig:section7_fig5}. Thus, this modeling approach can be used to address the data sparsity issues present in the early years of the augmented PRC dataset. 
%
Secondly, when calculating the month-to-month breach-to-victim conversion rate using the number of unique individuals rather than time-discounted cumulative records (Figure \ref{fig:saturation_conversion_rate}), the resulting victim conversion rate remains relatively stable, fluctuating around a 1:1 ratio. 
%
Finally, despite the de-duplication of the record pool, the statistical relationship between mega-breaches and IDT remains robust; the Wilcoxon signed-rank test results for the saturation model (Figure \ref{fig:wilcoxon_sat}) preserve the significance observed in the raw data, with $p$-values remaining below 0.05 for discovery lags of one and two months. Ultimately, the dynamic saturation model offers an alternative framework for measuring the ``yield'' of a breach in an increasingly crowded data market, and future work can include refining this model to account for the varying sensitivity of specific data types, such as SSNs, to better quantify long-term social harm.

% \subsection{Additional Future Work}
% \resp{TODO: add more future work here}
% \rev{Another line of future work involves improving the imputation process for breaches with undisclosed record counts. An alternative to our procedure (detailed in Section \ref{sec:imputation_of_undisclosed}) is to calculate a category-specific annual baseline $n_u^j$ for each of the $k$ breach types. Under this more granular approach, the imputed records for an incident $i$ of type $j$ in year $y$ would be defined as $R_{i_j}=n_u^j/|I_{u,y}^j|$. While this would account for the different typical magnitudes associated with various breach types (e.g. malicious attacks versus accidental disclosures), we maintained the global $n_u$ approach to provide a stable, conservative baseline for the overall number of undisclosed records. }
